% =========================================================
\section{Related Work}
\label{sec:related_work}
% =========================================================

This section reviews prior work on Koopman modeling and recursive prediction, temporal sequence modeling, and real-time predictive control.

\subsection{Koopman Learning and Recursive Prediction}

Koopman operator theory represents nonlinear dynamics as a linear evolution on the observable functions of the state \cite{koopman1931hamiltonian, koopman1932spectra, mezic2005spectral}, enabling the direct application of linear control synthesis to nonlinear systems. Extensions such as DMDc \cite{proctor2016dmdc}, input-aware formulations \cite{proctor2018generalizing, brunton2016invariant}, and Koopman-based MPC \cite{korda2018linear} established the foundation for Koopman control, and later work added learned representations and performance guarantees \cite{han2020deep, shi2022deep, strasser2026overview, duran2025control}. Controller performance in this literature is highly sensitive to model quality, which calls for improved lifting design.

To learn Koopman models from data, EDMD fits observables from a pre-defined dictionary \cite{williams2015data, williams2015kernel, klus2018data}, whereas data-driven methods learn the representations instead \cite{takeishi2017koopman, lusch2018deep, yeung2019learning, morton2019deep, yuan2022variational}. Most construct observables from instantaneous states and identify a single time-invariant operator, neglecting short-term temporal correlations and limiting long-horizon prediction accuracy.

Two lines relax this restriction. Delay-embedding variants enrich the dictionary with time-shifted measurements or Hankel/HAVOK coordinates \cite{kamb2020timedelay, brunton2017chaos, arbabi2017ergodic, das2019delay}, so the operator acts on a history-augmented state, but history enters only through the observables. Context-aware and time-varying formulations instead let the effective operator depend on recent data: Koopa computes context-aware operators for non-stationary forecasting \cite{liu2023koopa}, and MamKO generates time-varying Koopman matrices from online history for control \cite{li2025mamko}. Our formulation also lets history influence the transition, but confines it to a single low-rank family identified jointly with the lifting map and linear in the lifted state, so predictive control remains a quadratic program.

Multi-step training is also used in learned Koopman models to mitigate error accumulation under recursive prediction \cite{lusch2018deep, otto2019linearly, azencot2020forecasting, wu2026multistep}, and multi-step simulation-error objectives are standard in neural system identification more broadly \cite{forgione2021dynonet}; we therefore do not claim multi-step loss itself as novel. The issue here is specific to history-conditioned dynamics: during recursive deployment the transition is repeatedly rescheduled from the predicted history, so the propagated latent state and conditioning context can depart from the encoder-consistent one-step distribution. Our rollout-aware procedure reproduces that scheduled propagation during training rather than merely extending a one-step objective over several horizons.

\subsection{Transformers for Temporal Sequence Modeling}

Transformers \cite{vaswani2017attention} model temporal dependencies via self-attention, capturing dependencies across all time steps simultaneously and parallelizing better than step-by-step recurrent networks. Efficient variants and patch-based tokenization \cite{zhou2021informer, wang2020linformer, choromanski2021performer, wu2021autoformer, zhou2022fedformer, nie2023time} further improve efficiency by aggregating local temporal structure before global attention. In Koopman lifting, where the relevant temporal context is not known a priori, Transformers can incorporate short-term information from past trajectories.

Recent work has explored Transformer encoders with Koopman-style latent dynamics for forecasting \cite{forootani2026deepkoopformer, forootani2026spectral}. In particular, \cite{forootani2026spectral} parameterizes the transition as a single low-rank operator with a spectrally constrained diagonal factor on a patch-based Transformer backbone. It shares our patch-based encoder and low-rank factorization but differs in two respects that matter for control: its operator is a static low-rank transition used for forecasting, whereas we keep a full-rank nominal transition and let history modulate only a low-rank deviation; and our prediction is generated by the identified Koopman model inside a real-time predictive controller rather than by the Transformer itself. Relatedly, low-rank space-varying transitions have been learned for continuum-robot control \cite{feizi2025deepkoopman}, where auxiliary networks parameterize the state dependence of the transition spectrum and input matrix; that formulation is conditioned on the instantaneous state and paired with an infinite-horizon linear-quadratic controller rather than a constrained predictive one.

\subsection{Real-Time and Scheduled Predictive Control}

Two established strategies reduce the online cost of predictive control. The real-time iteration (RTI) performs a single sequential quadratic programming step per sampling instant around the current nominal trajectory, replacing the nonlinear program by a quadratic program \cite{diehl2005realtime, gros2020rti}; its nominal stability and suboptimality have been analyzed \cite{diehl2005rtistability}, and sensitivity-based schemes reduce computational delay by separating an offline nonlinear program from a fast online update \cite{zavala2009advancedstep}. Scheduled, or quasi-linear parameter-varying (quasi-LPV), predictive control instead embeds a nonlinear model into a family of linear models indexed by a scheduling variable, so that freezing or iterating the scheduling trajectory turns the problem into a sequence of QPs \cite{rugh2000gain, leith1998velocity, cisneros2016qlmpc, cisneros2020qlpv}.

Our controller inherits the one-iteration RTI structure but differs in where the scheduling signal comes from. In quasi-LPV MPC the scheduling variable is an explicit function of the state or velocity, with its predicted trajectory taken from the previous solution or a planned reference. Here the transition is scheduled by a context that a learned encoder extracts from a history window, and during prediction the history is itself produced by the predictor, so the scheduling trajectory is endogenous to the learned model and cannot be read off the measured state. Consequently, the context must be re-evaluated along the rolled-out latent trajectory rather than assumed, while the encoder need not be re-run at every prediction stage, letting the single-iteration controller fit inside the sampling period.
